% QUESTAO
Dadas duas matrizes $A$ e $B$ de mesma ordem, denominamos diferença entre $A$ e $B$, indicada por $A-B$, a matriz $C$ obtida ao calcularmos a adição de $A$ com o oposto de $B$, ou seja, $A-B=A+(-B)=C$. Assim, dada as matriz $A = \left[ \begin{array}{cc}
1 & 0 \\
-2 & 4
\end{array} \right]$ e $B = \left[ \begin{array}{cc}
8 & -2 \\
-5 & 1
\end{array} \right]$, determine a diferença $A-B$.

% ALTERNATIVAS
$A-B=\left[ \begin{array}{cc}	-7 & 2 \\	3 & 3 \end{array} \right]$
$A-B=\left[ \begin{array}{cc}	7 & 2 \\	3 & 5 \end{array} \right]$
$A-B=\left[ \begin{array}{cc}	3 & -2 \\	3 & 3 \end{array} \right]$
$A-B=\left[ \begin{array}{cc}	-7 & 0 \\	5 & 3 \end{array} \right]$
$A-B=\left[ \begin{array}{cc}	7 & 6 \\	1 & 3 \end{array} \right]$


